\apendice{Anexo de sostenibilización curricular}

\section{Introducción}

Este anexo tiene como objetivo reflexionar sobre los aspectos de sostenibilidad abordados durante la realización del presente Trabajo Fin de Grado y sobre las competencias sobre la sostenibilidad que han sido aplicadas. Para ello, se han tomado como referencia las Directrices para la introducción de la sostenibilidad en el currículum elaboradas por la CRUE~\cite{ODS:CRUE} y los Objetivos de Desarrollo Sostenible (ODS)~\cite{ODS:Naciones_unidas} establecidos en la Agenda 2030 de las Naciones Unidas.

Aunque PythIA es un proyecto de carácter tecnológico, durante su desarrollo se han abordado aspectos relacionados con la accesibilidad a la información, la transparencia administrativa, la eficiencia en el uso de recursos y el empleo responsable de sistemas de inteligencia artificial.

\section{ODS 9: Industria, innovación e infraestructura}

El ODS 9 persigue promover la innovación y desarrollar infraestructuras fiables, sostenibles y accesibles. Este objetivo guarda una relación directa con PythIA, ya que la aplicación utiliza tecnologías de inteligencia artificial para facilitar el acceso y la explotación de la información contenida en grandes volúmenes de documentación administrativa y técnica.

La principal aportación del proyecto consiste en mejorar la forma en que los usuarios acceden al conocimiento almacenado en los documentos. Frente a los métodos tradicionales de búsqueda manual, PythIA permite realizar consultas en lenguaje natural y obtener de forma rápida la información relevante, contribuyendo a modernizar y hacer más eficientes los procesos de gestión documental.

Además, durante el desarrollo del proyecto se ha optado por emplear tecnologías abiertas y ampliamente utilizadas en el ámbito profesional y científico. El uso de herramientas de código abierto favorece la reutilización del conocimiento, facilita la colaboración entre organizaciones y reduce la dependencia de soluciones propietarias. Este enfoque contribuye a la construcción de infraestructuras digitales más accesibles, sostenibles y adaptables a las necesidades futuras.

Desde la perspectiva de la sostenibilidad, el proyecto pone de manifiesto cómo la innovación tecnológica puede utilizarse para mejorar procesos existentes, incrementar la eficiencia en el acceso a la información y generar soluciones que aporten valor tanto a las organizaciones como a la sociedad.

\section{ODS 12: Producción y consumo responsables}

El ODS 12 promueve un uso eficiente y responsable de los recursos. En este sentido, PythIA contribuye a optimizar la gestión y explotación de la información contenida en grandes colecciones documentales, reduciendo el esfuerzo necesario para localizar y analizar datos relevantes.

En numerosos procedimientos administrativos se deben revisar una gran cantidad de documentación para encontrar información específica. Este proceso requiere una importante inversión de tiempo y recursos. PythIA facilita estas tareas mediante mecanismos de búsqueda semántica e inteligencia artificial que permiten acceder de forma rápida a la información relevante, favoreciendo una utilización más eficiente del conocimiento disponible. La indexación automática de documentos y la reutilización de la información previa evitan duplicidades en las tareas de análisis documental. Esto permite aprovechar mejor los recursos disponibles y reducir trabajos repetitivos que aportan poco valor añadido.

Por otra parte, la aplicación incorpora mecanismos de trazabilidad que permiten identificar los fragmentos documentales utilizados para generar cada respuesta. Esta característica fomenta un uso responsable de la inteligencia artificial, ya que facilita la verificación de la información proporcionada y ayuda a que las decisiones se apoyen en evidencias documentales verificables.

\section{ODS 13: Acción por el clima}

Aunque PythIA no ha sido desarrollado específicamente para abordar problemas medioambientales, el proyecto incorpora algunas prácticas que contribuyen indirectamente a un uso más eficiente de los recursos.

La aplicación facilita la consulta digital de grandes volúmenes de documentación, reduciendo la necesidad de imprimir expedientes y favoreciendo el acceso electrónico a la información. Esto contribuye a disminuir el consumo de papel y simplifica la gestión documental.

Además, durante el desarrollo se han adoptado soluciones orientadas a mejorar la eficiencia del sistema, reutilizando la información previamente procesada y empleando herramientas especializadas para cada tarea. Estas decisiones permiten optimizar el uso de los recursos computacionales necesarios para el funcionamiento de la aplicación.

\section{ODS 16: Paz, justicia e instituciones sólidas}

Uno de los aspectos más relevantes de PythIA es su contribución a mejorar el acceso a la información pública relacionada con los procedimientos de contratación.

Los expedientes administrativos suelen contener una gran cantidad de documentación técnica y jurídica cuya consulta puede resultar compleja y requerir una inversión considerable de tiempo. PythIA facilita la búsqueda y localización de información mediante consultas en lenguaje natural, permitiendo acceder de forma más rápida y sencilla al contenido de estos documentos.

Esta mejora en la accesibilidad favorece la transparencia y contribuye a reducir las barreras de acceso a la información pública. Además, la aplicación incorpora mecanismos de trazabilidad que permiten consultar los fragmentos documentales utilizados para generar cada respuesta, facilitando la comprobación de la información obtenida.

\section{Competencias de sostenibilidad adquiridas}

La realización de este Trabajo de Fin de Grado me ha permitido comprender que el desarrollo de soluciones informáticas implica considerar aspectos que van más allá de los requisitos puramente técnicos. Durante el desarrollo de PythIA he podido reflexionar sobre cómo una aplicación basada en inteligencia artificial puede influir en la forma en que las personas acceden a la información y toman decisiones.

El proyecto me ha ayudado a adoptar una visión más global del problema abordado. El desarrollo de un sistema para consultar documentación pública no solo requiere diseñar algoritmos y desarrollar \eng{software}, sino también tener en cuenta cuestiones relacionadas con la transparencia, la accesibilidad de la información, la protección de datos y la fiabilidad de los resultados proporcionados por la inteligencia artificial.

Otra de las competencias adquiridas ha sido la capacidad de analizar de forma crítica las tecnologías empleadas. Durante el desarrollo fue necesario evaluar diferentes modelos de lenguaje, sistemas de recuperación de información y mecanismos de evaluación, valorando tanto su rendimiento como su capacidad para proporcionar respuestas fiables y verificables.

En conjunto, este trabajo me ha permitido comprobar que la sostenibilidad también puede abordarse desde la ingeniería informática. Aspectos como la accesibilidad a la información, la transparencia, la eficiencia en el uso de los recursos y el desarrollo responsable de sistemas basados en inteligencia artificial forman parte de los retos que deben considerarse durante el diseño de soluciones tecnológicas. Esta experiencia ha reforzado la importancia de integrar criterios éticos, sociales y de sostenibilidad en el ejercicio de la profesión.